% A math character whose family has never been given a font.
%
% Measured against pdftex: named and dropped --
%
%   ! \textfont 0 is undefined (character ^^G).
%   Somewhere in the math formula just ended, you used the
%   stated character from an undefined font family. For example,
%   plain TeX doesn't allow \it or \sl in subscripts. Proceed,
%   and I'll try to forget that I needed that character.
%
% The size is named as well as the family, so the same character in a
% subscript draws \scriptfont and one two levels down \scriptscriptfont.
% The character is written the way the log writes any character.
%
% The box is the same either way -- the character was already being
% dropped here -- so this is a difference in what is reported and not in
% what is set.
%
% The ligature pass looks characters up too, to decide about kerns. It
% says nothing: the reading that wants the character is the one that
% reports for it, or `$\mathchar"07\mathchar"08$' would name the first of
% them twice.
%
% NOT YET, and it is not this message's doing: the reference writes
% character 200 as `^^c8' in trip's log and as the byte itself in a run
% under the ordinary format. That is INITEX having eight-bit printing off
% where a preloaded format has it on -- it applies to everything printed,
% not to this message, and docs/DECISIONS.md's
% a-character-the-log-writes-as-itself was measured under the format that
% has it on.
\tracingonline=1\showboxbreadth=9999 \showboxdepth=9999
\textfont0=\nullfont \setbox0=\hbox{$\mathchar"0007$}\showbox0
\setbox1=\hbox{$\mathchar"0007\mathchar"0008$}\showbox1
{\scriptfont0=\nullfont \setbox2=\hbox{$x^{\mathchar"0007}$}\showbox2}
{\scriptscriptfont0=\nullfont \setbox3=\hbox{$x^{y^{\mathchar"0007}}$}\showbox3}
{\textfont5=\nullfont \setbox4=\hbox{$\mathchar"0507$}\showbox4}
\setbox5=\hbox{$\mathchar"0041$}\showbox5
\end
